% Options for packages loaded elsewhere
\PassOptionsToPackage{unicode}{hyperref}
\PassOptionsToPackage{hyphens}{url}
\PassOptionsToPackage{dvipsnames,svgnames,x11names}{xcolor}
\documentclass[
  a4paper,
]{ltjsarticle}
\usepackage{xcolor}
\usepackage[margin=20mm]{geometry}
\usepackage{amsmath,amssymb}
\setcounter{secnumdepth}{-\maxdimen} % remove section numbering
\usepackage{iftex}
\ifPDFTeX
  \usepackage[T1]{fontenc}
  \usepackage[utf8]{inputenc}
  \usepackage{textcomp} % provide euro and other symbols
\else % if luatex or xetex
  \usepackage{unicode-math} % this also loads fontspec
  \defaultfontfeatures{Scale=MatchLowercase}
  \defaultfontfeatures[\rmfamily]{Ligatures=TeX,Scale=1}
\fi
\usepackage{lmodern}
\ifPDFTeX\else
  % xetex/luatex font selection
  \setmainfont[]{DejaVu Serif}
  \setsansfont[]{DejaVu Sans}
  \setmonofont[]{DejaVu Sans Mono}
\fi
% Use upquote if available, for straight quotes in verbatim environments
\IfFileExists{upquote.sty}{\usepackage{upquote}}{}
\IfFileExists{microtype.sty}{% use microtype if available
  \usepackage[]{microtype}
  \UseMicrotypeSet[protrusion]{basicmath} % disable protrusion for tt fonts
}{}
\makeatletter
\@ifundefined{KOMAClassName}{% if non-KOMA class
  \IfFileExists{parskip.sty}{%
    \usepackage{parskip}
  }{% else
    \setlength{\parindent}{0pt}
    \setlength{\parskip}{6pt plus 2pt minus 1pt}}
}{% if KOMA class
  \KOMAoptions{parskip=half}}
\makeatother
\usepackage{longtable,booktabs,array}
\newcounter{none} % for unnumbered tables
\usepackage{calc} % for calculating minipage widths
% Correct order of tables after \paragraph or \subparagraph
\usepackage{etoolbox}
\makeatletter
\patchcmd\longtable{\par}{\if@noskipsec\mbox{}\fi\par}{}{}
\makeatother
% Allow footnotes in longtable head/foot
\IfFileExists{footnotehyper.sty}{\usepackage{footnotehyper}}{\usepackage{footnote}}
\makesavenoteenv{longtable}
\usepackage{graphicx}
\makeatletter
\newsavebox\pandoc@box
\newcommand*\pandocbounded[1]{% scales image to fit in text height/width
  \sbox\pandoc@box{#1}%
  \Gscale@div\@tempa{\textheight}{\dimexpr\ht\pandoc@box+\dp\pandoc@box\relax}%
  \Gscale@div\@tempb{\linewidth}{\wd\pandoc@box}%
  \ifdim\@tempb\p@<\@tempa\p@\let\@tempa\@tempb\fi% select the smaller of both
  \ifdim\@tempa\p@<\p@\scalebox{\@tempa}{\usebox\pandoc@box}%
  \else\usebox{\pandoc@box}%
  \fi%
}
% Set default figure placement to htbp
\def\fps@figure{htbp}
\makeatother
\setlength{\emergencystretch}{3em} % prevent overfull lines
\providecommand{\tightlist}{%
  \setlength{\itemsep}{0pt}\setlength{\parskip}{0pt}}
\usepackage{xurl}
\usepackage{float}
\usepackage{bookmark}
\IfFileExists{xurl.sty}{\usepackage{xurl}
\usepackage{float}}{} % add URL line breaks if available
\urlstyle{same}
\hypersetup{
  pdftitle={三次元正方配列の三角形ループ積における例外ジョルダン代数 J\_3(\textbackslash mathbb\{O\}) 結晶、3 世代トライアリティ、および素の結合定数幾何統一（Part VII）},
  pdfauthor={木原 範昭（株式会社WFシステム）　ORCID 0009-0004-6753-4020},
  colorlinks=true,
  linkcolor={Maroon},
  filecolor={Maroon},
  citecolor={Blue},
  urlcolor={Blue},
  pdfcreator={LaTeX via pandoc}}

\title{三次元正方配列の三角形ループ積における例外ジョルダン代数
\(J_3(\mathbb{O})\) 結晶、3
世代トライアリティ、および素の結合定数幾何統一（Part VII）}
\usepackage{etoolbox}
\makeatletter
\providecommand{\subtitle}[1]{% add subtitle to \maketitle
  \apptocmd{\@title}{\par {\large #1 \par}}{}{}
}
\makeatother
\subtitle{― 非対角 3 世代巡回積 \(J_{\mathrm{off}}\)（24D）の
\(\mathrm{Tri}(T)\)
厳密不動点、\(F_4 \supset SO(8) \supset G_2 \supset SU(3)_C\) 連鎖による
3 世代同定、非線形階層増幅、および素の統一比
\(\alpha_3:\alpha_2:\alpha_1 = 1:1:1\) の監査 ―}
\author{木原 範昭（株式会社WFシステム）　ORCID 0009-0004-6753-4020}
\date{2026年11月28日　Version DOI 10.5281/zenodo.24628777}

\begin{document}
\maketitle

{
\setcounter{tocdepth}{2}
\tableofcontents
}
\textbf{著者：} 木原 範昭（株式会社WFシステム）\\
\textbf{ORCID：} 0009-0004-6753-4020\\
\textbf{日付：} 2026年11月28日\\
\textbf{Version DOI：} 10.5281/zenodo.24628777\\
\textbf{Concept DOI：} 10.5281/zenodo.24628776\\
\textbf{位置づけ：} 三角形ループ積シリーズ Part I--VI（Concept DOIs:
\texttt{24622200}, \texttt{24622565}, \texttt{24625962},
\texttt{24626592}, \texttt{24627519},
\texttt{24628297}）の続編となる\textbf{3
世代構造・結合定数統一・質量階層監査論文（Part VII）}。これまで 1
世代（\(O_8 \otimes O_8\), 64D）を扱ってきた本シリーズに対し、八元数 3
個を包含する例外ジョルダン代数 \(J_3(\mathbb{O})\)（27D）を
\(\mathrm{Tri}(T)\) 動力学へ直接投入し、(1) \(J_3(\mathbb{O})\)
のどの部分が \(\mathrm{Tri}(T)\) の厳密不動点結晶になるか、(2)
\(F_4 \supset SO(8) \supset G_2 \supset SU(3)_C\) による 3
世代トライアリティの機械精度証明、(3) 対称性破れによる 3
世代階層増幅と小出比 \(Q\) の反証監査、(4)
無パラメータの素の結合定数統一比と低エネルギー実測値との RG
境界を、決定論的検証スクリプト（\texttt{tri\_loop\_part7\_jordan\_generations\_v0.py}、乱数なし）により白黒つける。

\begin{center}\rule{0.5\linewidth}{0.5pt}\end{center}

\subsection{要旨}\label{ux8981ux65e8}

標準模型の最大の謎の一つである「なぜフェルミオンは 3
世代存在するのか」および「結合定数・質量階層はどこまで純粋代数・\(\mathrm{Tri}(T)\)
動力学だけで決まるのか」を検証するため、本 Part VII
では例外ジョルダン代数 \(J_3(\mathbb{O})\)（27D：対角
\(\mathbb{R}^3\)＋非対角 3 八元数
\(\mathbb{O}_1 \oplus \mathbb{O}_2 \oplus \mathbb{O}_3\)）を構成し、3
次元正方配列の三角形ループ積
\(\mathrm{Tri}(T)_{ijk} = \sum_{p,q,r} T_{ipq} T_{qjr} T_{rpk}\)
との厳密接続を機械精度で検証した。

\begin{enumerate}
\def\labelenumi{\arabic{enumi}.}
\tightlist
\item
  \textbf{非対角 3 世代巡回積 \(J_{\mathrm{off}}\)（24D）の
  \(\mathrm{Tri}(T)\) 厳密結晶定理 {[}J3-1{]}}：ジョルダン恒等式誤差
  \texttt{8.269e-15} で構成した \(J_3(\mathbb{O})\) において、対角 3
  成分 \(\{E_1,E_2,E_3\}\) を含むフル積テンソル \(J\) や 3 次不変式
  \(D\) は \(\mathrm{Tri}(T)\) の反復で rank-1
  対角点（\(\|P\| = 1.000000\)）へ崩壊する一方、\textbf{非対角 3
  八元数ブロック間の巡回積
  \(J_{\mathrm{off}}\)（24D：\(\mathrm{Re}((xy)z)\)）は
  \(\mathrm{Tri}(J_{\mathrm{off}}) = -\frac{1}{64} J_{\mathrm{off}}\)
  を誤差 \texttt{0.000e+00} で満たす厳密不動点結晶}（100
  ステップ後も誤差 \texttt{0.000e+00}、非対角摂動に対し
  \texttt{3.536e-05} で自己修復）であることを発見した。
\item
  \textbf{\(F_4(52\mathrm{D}) \supset SO(8)(28\mathrm{D}) \supset G_2(14\mathrm{D}) \supset SU(3)_C(8\mathrm{D})\)
  による 3 世代トライアリティの機械精度証明 {[}J3-2{]}}：\(J_{ijk}\)
  の導分代数 SVD
  核空間から、\(\mathrm{Der}(J_3(\mathbb{O})) = \mathfrak{f}_4\)（\textbf{厳密
  52D}、最大ゼロ特異値 \texttt{1.996e-15}）、対角固定部分代数
  \(\mathrm{Der}_0 = \mathfrak{so}(8)\)（\textbf{厳密 28D}、世代間混合
  \texttt{1.285e-15}）を抽出。第 1 世代 \(O_1\) の \(e_0\) を固定する
  \(\mathfrak{so}(7)_1\)（\textbf{21D}）の各世代固定ベクトル次元は
  \textbf{\texttt{{[}1,\ 0,\ 0{]}}} となり、3
  世代ブロックが互いに非同値なトライアリティ表現
  \textbf{\(8_v \oplus 8_s \oplus 8_c\)} であることを証明した。さらに全
  3 世代同時固定群 \(\mathfrak{g}_2\)（\textbf{14D}）および
  \(SU(3)_C\)（\textbf{8D}）の下では、全 3 世代の二次カシミール固有値が
  \textbf{\texttt{{[}0,\ 0,\ 8/9,\ 8/9,\ 8/9,\ 8/9,\ 8/9,\ 8/9{]}}}（レプトン二重項
  \texttt{2}＋クォーク \(3\oplus\bar{3}\) \texttt{6}）で完全一致する。
\item
  \textbf{3 世代非線形階層増幅と小出比 \(Q\) の反証境界
  {[}J3-3{]}}：厳密結晶 \(J_{\mathrm{off}}\)（固有値
  \texttt{1/64}）に対し、各世代の局在自己結合摂動（固有値 \texttt{1}、64
  倍の増幅率）を 2 段階で投入すると、微小な 2 段階差から \(\tau = 40\)
  で
  \texttt{E(O1):E(O2):E(O3)\ =\ 500.4\ :\ 17.2\ :\ 1.0}（\texttt{(0.14,\ 0.08)}）や
  \texttt{2138.6\ :\ 16.0\ :\ 1.0}（\texttt{(0.18,\ 0.12)}）という
  \textbf{2〜3 桁の指数的 3 世代階層} が自発増幅される。一方、小出比
  \(Q(E)\) は
  \texttt{0.333\ \$\textbackslash{}to\$\ 0.685\ \$\textbackslash{}to\$\ 0.731\ \$\textbackslash{}to\$\ 0.821\ \$\textbackslash{}to\$\ 1.000}
  と初期摂動比に連続依存し、普遍定数 \(Q = 2/3\)
  へは自動ロックしないことを反証境界として明示した。
\item
  \textbf{素の結合定数幾何統一比
  \(\alpha_3^{\mathrm{bare}} : \alpha_2^{\mathrm{bare}} : \alpha_1^{\mathrm{bare}} = 1:1:1\)
  と RG 境界 {[}J3-4{]}}：64D 1 世代
  \(O_8^{\mathrm{ext}}\otimes O_8^{\mathrm{int}}\)
  上の標準模型生成子トレースは
  \(\mathrm{Tr}_{64}(T_3^{C\,2}) = \mathrm{Tr}_{64}(T_3^{L\,2}) = \mathrm{Tr}_{64}(Y_1^2) = \mathbf{8.000000}\)（GUT
  正規化
  \(Y_1 = \sqrt{3/5}\,Y\)）、\(\mathrm{Tr}_{64}(Q_{\mathrm{em}}^2) = 64/3\)
  となり、調整パラメータゼロで
  \textbf{\(\alpha_3^{\mathrm{bare}} : \alpha_2^{\mathrm{bare}} : \alpha_1^{\mathrm{bare}} = 1 : 1 : 1\)}、\textbf{\(\sin^2\theta_W^{\mathrm{bare}} = \alpha_{\mathrm{em}}^{\mathrm{bare}}/\alpha_s^{\mathrm{bare}} = 3/8 = 0.375000\)}
  が導出される（27D \(J_3(\mathbb{O})\) の 3 世代 \(SU(3)_C\) トレースも
  \texttt{1:1:1}
  完全均等）。低エネルギー実測値（\(\alpha \approx 1/137.036, \alpha_s(M_Z) \approx 0.1179, \sin^2\theta_W(M_Z) \approx 0.2312\)）との差は数合わせを行わず、繰り込み群（RG）走行の未検証境界として明示的に切り分ける。
\end{enumerate}

\begin{center}\rule{0.5\linewidth}{0.5pt}\end{center}

\subsection{\texorpdfstring{1. 導入：1 世代 \(O_8 \otimes O_8\) から 3
世代例外ジョルダン代数 \(J_3(\mathbb{O})\)
へ}{1. 導入：1 世代 O\_8 \textbackslash otimes O\_8 から 3 世代例外ジョルダン代数 J\_3(\textbackslash mathbb\{O\}) へ}}\label{ux5c0eux51651-ux4e16ux4ee3-o_8-otimes-o_8-ux304bux3089-3-ux4e16ux4ee3ux4f8bux5916ux30b8ux30e7ux30ebux30c0ux30f3ux4ee3ux6570-j_3mathbbo-ux3078}

Part I--VI {[}P1--P6{]} は、単一八元数ペア
\(O_8 \otimes O_8\)（64D）から標準模型 1
世代のゲージ群・パリティ最大破れ・多殻重力・2 粒子散乱・CHSH
ベル不等式破れを実証した。一方、素粒子標準模型の 3
世代構造を八元数から自然に導く数学的構造として、3×3
エルミート八元数行列のなす 27 次元例外ジョルダン代数
\(J_3(\mathbb{O})\)（Albert 代数）{[}E1--E3{]}
が古くから注目されている。

本 Part VII の核心は、\textbf{「\(J_3(\mathbb{O})\) は三角形ループ積
\(\mathrm{Tri}(T)\) の動力学の中で実際に結晶として選ばれるのか？」}
を直接テストし、3
世代同定・質量階層・素の結合定数統一比の到達点と限界を、\textbf{【既知】・【Part
I--VI 確立】・【Part VII 検証】・【仮説・未検証】} の 4
段階で峻別することである。

\begin{center}\rule{0.5\linewidth}{0.5pt}\end{center}

\subsection{\texorpdfstring{2. 例外ジョルダン代数 \(J_3(\mathbb{O})\) の
\(\mathrm{Tri}(T)\) 結晶性【Part VII 検証
J3-1】}{2. 例外ジョルダン代数 J\_3(\textbackslash mathbb\{O\}) の \textbackslash mathrm\{Tri\}(T) 結晶性【Part VII 検証 J3-1】}}\label{ux4f8bux5916ux30b8ux30e7ux30ebux30c0ux30f3ux4ee3ux6570-j_3mathbbo-ux306e-mathrmtrit-ux7d50ux6676ux6027part-vii-ux691cux8a3c-j3-1}

\(J_3(\mathbb{O})\) の元
\(X = \begin{pmatrix} a_1 & z & \bar{y} \\ \bar{z} & a_2 & x \\ y & \bar{x} & a_3 \end{pmatrix}\)（\(a_i \in \mathbb{R},\; x, y, z \in \mathbb{O}\)）に対し、正規直交内積
\(\langle X, Y \rangle = \mathrm{Tr}(X \circ Y)\) のもとで 27
次元基底を構成した： - 対角基底（3D）：\(E_1, E_2, E_3\)（指標
\texttt{0,\ 1,\ 2}） - 非対角 3
世代基底（24D）：\(O_1\)（\texttt{3..10}）、\(O_2\)（\texttt{11..18}）、\(O_3\)（\texttt{19..26}）

ジョルダン積 \(X \circ Y = \frac{1}{2}(XY+YX)\) の構造定数テンソル
\(J_{ijk}\) はジョルダン恒等式
\(|(X^2\circ Y)\circ X - X^2\circ(Y\circ X)| \le \mathbf{8.269 \times 10^{-15}}\)・完全対称性誤差
\texttt{1.110e-16} を満たす。これを 3 種（フル \(J\)、3
次不変式行列式テンソル \(D\)、および非対角 3 世代巡回積
\(J_{\mathrm{off}}\)）に分けて \(\mathrm{Tri}(T)\) へ投入した結果を表 1
に示す：

{\def\LTcaptype{none} % do not increment counter
\begin{longtable}[]{@{}
  >{\raggedright\arraybackslash}p{(\linewidth - 10\tabcolsep) * \real{0.1379}}
  >{\centering\arraybackslash}p{(\linewidth - 10\tabcolsep) * \real{0.1724}}
  >{\centering\arraybackslash}p{(\linewidth - 10\tabcolsep) * \real{0.1724}}
  >{\centering\arraybackslash}p{(\linewidth - 10\tabcolsep) * \real{0.1724}}
  >{\centering\arraybackslash}p{(\linewidth - 10\tabcolsep) * \real{0.1724}}
  >{\centering\arraybackslash}p{(\linewidth - 10\tabcolsep) * \real{0.1724}}@{}}
\toprule\noalign{}
\begin{minipage}[b]{\linewidth}\raggedright
投入テンソル
\end{minipage} & \begin{minipage}[b]{\linewidth}\centering
次元
\end{minipage} & \begin{minipage}[b]{\linewidth}\centering
1-step \(\|\mathrm{Tri}(T_0)\|\)
\end{minipage} & \begin{minipage}[b]{\linewidth}\centering
1-step 不動点誤差 \(\|\pm T_1 - T_0\|\)
\end{minipage} & \begin{minipage}[b]{\linewidth}\centering
100-step 不動点誤差
\end{minipage} & \begin{minipage}[b]{\linewidth}\centering
判定
\end{minipage} \\
\midrule\noalign{}
\endhead
\bottomrule\noalign{}
\endlastfoot
フルジョルダン積 \(J\) & 27D & \texttt{0.019909} (\texttt{+}) &
\texttt{9.617e-01} & \texttt{1.370e+00} (\(\|P\|=1.0\)) & 対角 rank-1
崩壊 \\
3 次不変式 \(D\) (\(\det\)) & 27D & \texttt{0.014784} (\texttt{-}) &
\texttt{1.037e+00} & \texttt{1.364e+00} (\(\|P\|=1.0\)) & 対角 rank-1
崩壊 \\
\textbf{非対角 3 世代積 \(J_{\mathrm{off}}\)} & \textbf{24D} &
\textbf{\texttt{0.015625} (\texttt{-1/64})} &
\textbf{\texttt{0.000e+00}} & \textbf{\texttt{0.000e+00}} &
\textbf{厳密不動点結晶} \\
\end{longtable}
}

\textbf{定理的発見}：\(J_3(\mathbb{O})\) の 3
つの非対角八元数ブロック間の巡回積
\(J_{\mathrm{off}}\)（\(\mathrm{Re}((xy)z)\)）は、まさに 3 指標
\((i,j,k)\) が 3 世代 \((O_1, O_2, O_3)\)
を巡回する三角形ループ構造そのものであり、\(\mathrm{Tri}(J_{\mathrm{off}}) = -\frac{1}{64} J_{\mathrm{off}}\)
の\textbf{誤差ゼロ厳密不動点結晶}をなす（非対角微小摂動 \texttt{1e-4}
に対しても 100 ステップ後 \texttt{3.536e-05} へ収束）。

\begin{center}\rule{0.5\linewidth}{0.5pt}\end{center}

\subsection{\texorpdfstring{3.
\(F_4 \supset SO(8) \supset G_2 \supset SU(3)_C\) 連鎖と 3
世代トライアリティ【Part VII 検証
J3-2】}{3. F\_4 \textbackslash supset SO(8) \textbackslash supset G\_2 \textbackslash supset SU(3)\_C 連鎖と 3 世代トライアリティ【Part VII 検証 J3-2】}}\label{f_4-supset-so8-supset-g_2-supset-su3_c-ux9023ux9396ux3068-3-ux4e16ux4ee3ux30c8ux30e9ux30a4ux30a2ux30eaux30c6ux30a3part-vii-ux691cux8a3c-j3-2}

\(J_{ijk}\) の導分条件
\(M_{mi} J_{mjk} + M_{mj} J_{imk} + M_{mk} J_{ijm} = 0\)（\(\mathfrak{so}(27)\)
の 351 基底、完全対称独立 3,654 成分）を SVD で厳密求解した（図 1(a)）：
1.
\textbf{\(\mathrm{Der}(J_3(\mathbb{O})) = \mathfrak{f}_4\)（52D）}：最小非ゼロ特異値
\texttt{1.732e+00}、最大ゼロ特異値 \textbf{\texttt{1.996e-15}}。 2.
\textbf{対角固定部分代数
\(\mathrm{Der}_0 = \mathfrak{so}(8)\)（28D）}：3 世代間の混合成分は最大
\textbf{\texttt{1.285e-15}}（厳密ブロック対角：3 世代が独立に閉じた 8D
表現をなす）。 3.
\textbf{トライアリティ非同値性（\(8_v \oplus 8_s \oplus 8_c\)）}：第 1
世代 \(O_1\) の実部 \(e_0\) を固定する部分代数
\(\mathfrak{so}(7)_1\)（\textbf{21D}）が各世代で固定するベクトル次元を計測すると
\textbf{\texttt{(O1,\ O2,\ O3)\ =\ {[}1,\ 0,\ 0{]}}} となった。これは
\(O_1\) がベクトル表現 \(8_v\)、\(O_2, O_3\)
が不変ベクトルを持たないスピノル表現 \(8_s, 8_c\)
であることの機械精度証明である。 4. \textbf{ゲージ群
\(G_2\)（14D）\(\supset SU(3)_C\)（8D）下の 3 世代同一コピー}：全 3
世代の \(e_0\) 同時固定群は \textbf{14D（\(G_2\)）}、さらに \(e_1\)
を固定すると \textbf{8D（\(SU(3)_C\)）} となる。この \(SU(3)_C\)
の二次カシミール固有値は、第 1・第 2・第 3 世代すべてで厳密に
\textbf{\texttt{{[}0,\ 0,\ 8/9,\ 8/9,\ 8/9,\ 8/9,\ 8/9,\ 8/9{]}}}（レプトン二重項
\texttt{2} ＋ カラー 3 重項・反 3 重項 \texttt{6}）に一致する。

\begin{center}\rule{0.5\linewidth}{0.5pt}\end{center}

\subsection{\texorpdfstring{4. 3 世代非線形階層増幅と小出比 \(Q\)
の反証境界【Part VII 検証
J3-3】}{4. 3 世代非線形階層増幅と小出比 Q の反証境界【Part VII 検証 J3-3】}}\label{ux4e16ux4ee3ux975eux7ddaux5f62ux968eux5c64ux5897ux5e45ux3068ux5c0fux51faux6bd4-q-ux306eux53cdux8a3cux5883ux754cpart-vii-ux691cux8a3c-j3-3}

厳密結晶 \(J_{\mathrm{off}}\)（固有値
\texttt{1/64\ =\ 0.015625}）に対し、各世代 \(O_g\) の局在自己結合摂動
\(P_{g} = \delta_{g,g,g}\)（固有値 \texttt{1.000}、\(J_{\mathrm{off}}\)
の \textbf{64 倍}）を投入し、\(\tau = 40\) での世代エネルギー占有率
\(E(O_1), E(O_2), E(O_3)\) と小出比
\(Q(E) = \sum E_g / (\sum \sqrt{E_g})^2\) を計測した（図 1(b)）：

{\def\LTcaptype{none} % do not increment counter
\begin{longtable}[]{@{}
  >{\centering\arraybackslash}p{(\linewidth - 10\tabcolsep) * \real{0.1667}}
  >{\centering\arraybackslash}p{(\linewidth - 10\tabcolsep) * \real{0.1667}}
  >{\centering\arraybackslash}p{(\linewidth - 10\tabcolsep) * \real{0.1667}}
  >{\centering\arraybackslash}p{(\linewidth - 10\tabcolsep) * \real{0.1667}}
  >{\centering\arraybackslash}p{(\linewidth - 10\tabcolsep) * \real{0.1667}}
  >{\centering\arraybackslash}p{(\linewidth - 10\tabcolsep) * \real{0.1667}}@{}}
\toprule\noalign{}
\begin{minipage}[b]{\linewidth}\centering
投入摂動 \((\epsilon_1, \epsilon_2)\)
\end{minipage} & \begin{minipage}[b]{\linewidth}\centering
\(E(O_1)\)
\end{minipage} & \begin{minipage}[b]{\linewidth}\centering
\(E(O_2)\)
\end{minipage} & \begin{minipage}[b]{\linewidth}\centering
\(E(O_3)\)
\end{minipage} & \begin{minipage}[b]{\linewidth}\centering
世代階層比 \(E_1:E_2:E_3\)
\end{minipage} & \begin{minipage}[b]{\linewidth}\centering
小出比 \(Q(E)\)
\end{minipage} \\
\midrule\noalign{}
\endhead
\bottomrule\noalign{}
\endlastfoot
\texttt{(0.00,\ 0.00)}（対称結晶） & \texttt{0.333333} &
\texttt{0.333333} & \texttt{0.333333} & \texttt{1.0\ :\ 1.0\ :\ 1.0} &
\texttt{0.33333} \\
\texttt{(0.010,\ 0.005)}（微小 2 段階） & \texttt{0.775423} &
\texttt{0.221966} & \texttt{0.002610} &
\textbf{\texttt{297.1\ :\ 85.0\ :\ 1.0}} & \texttt{0.50817} \\
\texttt{(0.14,\ 0.08)}（中規模 2 段階） & \texttt{0.964830} &
\texttt{0.033242} & \texttt{0.001928} &
\textbf{\texttt{500.4\ :\ 17.2\ :\ 1.0}} & \texttt{0.68472} \\
\texttt{(0.18,\ 0.12)}（強 2 段階） & \texttt{0.992123} &
\texttt{0.007413} & \texttt{0.000464} &
\textbf{\texttt{2138.6\ :\ 16.0\ :\ 1.0}} & \texttt{0.82093} \\
\end{longtable}
}

\begin{itemize}
\tightlist
\item
  \textbf{確立された機構}：固有値比 \texttt{64:1}
  の非線形不安定性により、わずか 2
  倍程度の初期対称性破れ（\(\epsilon_1 : \epsilon_2 = 2:1\)）から、\(\mathrm{Tri}(T)\)
  反復だけで \textbf{2〜3 桁に及ぶ指数的な 3
  世代階層（\(m_3 \gg m_2 \gg m_1\) 型）} が自動増幅される。
\item
  \textbf{反証・境界の明示}：小出比 \(Q(E)\) は初期摂動比
  \((\epsilon_1, \epsilon_2)\) に応じて \texttt{0.333} から
  \texttt{1.000} まで連続変化し、実験的な荷電レプトン値
  \(Q = 2/3 \approx 0.66667\) へ普遍的にロックする機構は単体
  \(J_{\mathrm{off}}\) の \(\mathrm{Tri}(T)\)
  には含まれない。精密質量比の第一原理決定には、外部空間・電弱ヒッグス結合との連成および
  RG 走行の導出が必要である【仮説・未検証】。
\end{itemize}

\begin{figure}[H]
\centering
\includegraphics[width=0.96\linewidth,height=\textheight,keepaspectratio,alt={Part VII 検証 (1)：(a) J\_3(\textbackslash mathbb\{O\})（27D）から SVD 核空間で厳密抽出された対称群連鎖 F\_4(52\textbackslash mathrm\{D\}) \textbackslash supset SO(8)(28\textbackslash mathrm\{D\}) \textbackslash supset SO(7)\_1(21\textbackslash mathrm\{D\}) \textbackslash supset G\_2(14\textbackslash mathrm\{D\}) \textbackslash supset SU(3)\_C(8\textbackslash mathrm\{D\}) とトライアリティ 3 世代、(b) 厳密結晶 J\_\{\textbackslash mathrm\{off\}\}（24D）上の 2 段階破れによる 3 世代非線形階層増幅。}]{figs_ja/fig_p7_j3o_triality_and_hierarchy_ja.png}
\caption{Part VII 検証 (1)：(a) \(J_3(\mathbb{O})\)（27D）から SVD
核空間で厳密抽出された対称群連鎖
\(F_4(52\mathrm{D}) \supset SO(8)(28\mathrm{D}) \supset SO(7)_1(21\mathrm{D}) \supset G_2(14\mathrm{D}) \supset SU(3)_C(8\mathrm{D})\)
とトライアリティ 3 世代、(b) 厳密結晶 \(J_{\mathrm{off}}\)（24D）上の 2
段階破れによる 3 世代非線形階層増幅。}
\end{figure}

\begin{center}\rule{0.5\linewidth}{0.5pt}\end{center}

\subsection{5. 素の結合定数幾何統一比（1:1:1）と低エネルギー RG
境界【Part VII 検証
J3-4】}\label{ux7d20ux306eux7d50ux5408ux5b9aux6570ux5e7eux4f55ux7d71ux4e00ux6bd4111ux3068ux4f4eux30a8ux30cdux30ebux30aeux30fc-rg-ux5883ux754cpart-vii-ux691cux8a3c-j3-4}

Part IV {[}P4{]} の 64D 1 世代表現
\(O_8^{\mathrm{ext}} \otimes O_8^{\mathrm{int}}\) 上の標準模型全 12
生成子に対し、無パラメータの幾何トレースを算出した（図 2(a)）： -
\(\mathrm{Tr}_{64}(T_3^{C\,2}) = \mathbf{8.000000}\) -
\(\mathrm{Tr}_{64}(T_3^{L\,2}) = \mathbf{8.000000}\) -
\(\mathrm{Tr}_{64}(Y^2) = \mathbf{40/3} \implies\) GUT 正規化
\(Y_1 = \sqrt{3/5}\,Y\) で
\(\mathrm{Tr}_{64}(Y_1^2) = \mathbf{8.000000}\) -
\(\mathrm{Tr}_{64}(Q_{\mathrm{em}}^2) = \mathbf{64/3}\)

これより、調整パラメータゼロの素の幾何結合定数比として
\[\alpha_3^{\mathrm{bare}} : \alpha_2^{\mathrm{bare}} : \alpha_1^{\mathrm{bare}} = 1 : 1 : 1, \qquad \sin^2\theta_W^{\mathrm{bare}} = \frac{\alpha_{\mathrm{em}}^{\mathrm{bare}}}{\alpha_s^{\mathrm{bare}}} = \frac{3}{8} = 0.375000\]
が厳密に得られる。また 27D 例外ジョルダン代数 \(J_3(\mathbb{O})\) 上の 3
世代別 \(SU(3)_C\) トレースも
\texttt{0.666667\ :\ 0.666667\ :\ 0.666667\ =\ 1\ :\ 1\ :\ 1}
で完全均等（3 世代ゲージ普遍性）である。

一方、低エネルギー実測値（\(M_Z\)
スケール：\(\sin^2\theta_W \approx 0.2312\)、\(\alpha_{\mathrm{em}}/\alpha_s \approx 0.0663\)、\(\alpha(0) \approx 1/137.036\)）との差は図
2(b)
の通り明確であり、\textbf{本モデルの素の幾何値は高エネルギー統一スケールに対応し、低エネルギー実測値への到達には繰り込み群（RG）走行の第一原理導出が必須である}
ことが数合わせなしに確定した【仮説・未検証】。

\begin{figure}[H]
\centering
\includegraphics[width=0.96\linewidth,height=\textheight,keepaspectratio,alt={Part VII 検証 (2)：(a) 64D 表現上の標準模型生成子幾何トレースによる厳密ゲージ統一 \textbackslash alpha\_3:\textbackslash alpha\_2:\textbackslash alpha\_1 = 1:1:1（\textbackslash mathrm\{Tr\}=8）、(b) 無パラメータの素の幾何統一値（青）と M\_Z 低エネルギー実測値（緑）の対比および RG 走行境界の明示。}]{figs_ja/fig_p7_coupling_unification_and_rg_boundary_ja.png}
\caption{Part VII 検証 (2)：(a) 64D
表現上の標準模型生成子幾何トレースによる厳密ゲージ統一
\(\alpha_3:\alpha_2:\alpha_1 = 1:1:1\)（\(\mathrm{Tr}=8\)）、(b)
無パラメータの素の幾何統一値（青）と \(M_Z\)
低エネルギー実測値（緑）の対比および RG 走行境界の明示。}
\end{figure}

\begin{center}\rule{0.5\linewidth}{0.5pt}\end{center}

\subsection{6.
シリーズ総括と再現手順}\label{ux30b7ux30eaux30fcux30baux7dcfux62ecux3068ux518dux73feux624bux9806}

本 Part VII をもって、当初設定した 4
大課題（電弱破れ・パリティ最大破れ［Part
IV］、多殻大格子重力・因果伝播・スピン 2 \(g_{00}(d)\)［Part V］、2
粒子動的散乱・空間分離 CHSH 破れ［Part VI］、および \(J_3(\mathbb{O})\)
3 世代トライアリティ・結合定数幾何統一［Part
VII］）を一巡して決定論的検証を完了した。

\begin{itemize}
\tightlist
\item
  \textbf{次段階の未解決課題}：(1) 多殻 \(K \ge 3\)
  におけるアインシュタイン方程式写像・2
  粒子波束散乱断面積の距離依存性、(2) 3 世代
  \(J_{\mathrm{off}} \otimes O_8^{\mathrm{ext}}\) からの CKM/PMNS
  混合角・CP 破れ、(3)
  高エネルギー幾何統一値（\(1:1:1,\; 3/8\)）から低エネルギー実測値（\(\alpha \approx 1/137.036, \alpha_s \approx 0.1179, \sin^2\theta_W \approx 0.2312\)）および精密質量比への繰り込み群走行の第一原理導出。
\item
  \textbf{再現ファイル}：\texttt{tri\_loop\_part7\_jordan\_generations\_v0.py}（J3-1〜J3-4、乱数なし・実行時間約
  1.6
  秒）、\texttt{results/tri\_loop\_part7\_jordan\_generations\_v0.txt}、\texttt{plot\_part7\_figures\_v0.py}、\texttt{results/fig\_p7\_*\_\{ja,en\}.png}。
\end{itemize}

\begin{center}\rule{0.5\linewidth}{0.5pt}\end{center}

\subsection{参考文献}\label{ux53c2ux8003ux6587ux732e}

\subsubsection{A. 自己引用}\label{a.-ux81eaux5df1ux5f15ux7528}

\begin{itemize}
\tightlist
\item
  \textbf{{[}P1{]}} 木原範昭,
  \emph{三次元正方配列の三角形ループ積からの量子・ゲージ・重力・フェルミオン構造の自発創発（統合完成論文
  Part I）}, Zenodo (2026-11-28), Concept DOI:
  \texttt{10.5281/zenodo.24622200}.
\item
  \textbf{{[}P2{]}} 木原範昭,
  \emph{三次元正方配列の三角形ループ積における重力・電磁気・QCD・ディラックの統一接続仮説（Part
  II）}, Zenodo (2026-11-28), Concept DOI:
  \texttt{10.5281/zenodo.24622565}.
\item
  \textbf{{[}P3{]}} 木原範昭,
  \emph{三次元正方配列の三角形ループ積における時空・内部空間読出し接続仮説（Part
  III）}, Zenodo (2026-11-28), Concept DOI:
  \texttt{10.5281/zenodo.24625962}.
\item
  \textbf{{[}P4{]}} 木原範昭,
  \emph{三次元正方配列の三角形ループ積における電弱対称性破れとパリティ最大破れ（Part
  IV）}, Zenodo (2026-11-28), Concept DOI:
  \texttt{10.5281/zenodo.24626592}.
\item
  \textbf{{[}P5{]}} 木原範昭,
  \emph{三次元正方配列の三角形ループ積における多殻大格子重力、因果伝播、およびスピン
  2 による \(g_{00}(d)\) 解禁（Part V）}, Zenodo (2026-11-28), Concept
  DOI: \texttt{10.5281/zenodo.24627519}.
\item
  \textbf{{[}P6{]}} 木原範昭,
  \emph{三次元正方配列の三角形ループ積における 2
  粒子動的散乱、ゲージ媒介相互作用、および空間分離 CHSH
  ベル不等式の動的破れ（Part VI）}, Zenodo (2026-11-28), Concept DOI:
  \texttt{10.5281/zenodo.24628297}.
\end{itemize}

\subsubsection{B.
外部一次文献}\label{b.-ux5916ux90e8ux4e00ux6b21ux6587ux732e}

\begin{itemize}
\tightlist
\item
  \textbf{{[}E1{]}} P. Jordan, J. von Neumann, and E. Wigner, \emph{On
  an Algebraic Generalization of the Quantum Mechanical Formalism}, Ann.
  of Math. \textbf{35}, 29--64 (1934).
\item
  \textbf{{[}E2{]}} J. C. Baez, \emph{The Octonions}, Bull. Amer. Math.
  Soc. \textbf{39}, 145--205 (2002), arXiv:math/0105155.
\item
  \textbf{{[}E3{]}} T. Dray and C. A. Manogue, \emph{The Geometry of the
  Octonions}, World Scientific (2015).
\item
  \textbf{{[}E4{]}} Y. Koide, \emph{Fermion-Boson Two-Body Model of
  Quarks and Leptons and Cabibbo Mixing}, Lett. Nuovo Cimento
  \textbf{34}, 201--205 (1982).
\end{itemize}

\end{document}
